\documentclass[10pt,a4paper]{amsart}
\usepackage[margin=19mm]{geometry}
\usepackage{amsmath,amssymb,mathtools}
\usepackage[scr=rsfs,cal=euler]{mathalfa}
\usepackage{xfrac}
\usepackage[unicode,hidelinks,bookmarks=false]{hyperref}
\newcommand{\ie}{i.e.\ }
\newcommand{\cf}{cf.\ }
\newcommand{\resp}{respectively\ }
\newcommand{\Subsection}{\subsection}
\providecommand{\nobreakdash}{\nobreak\hbox{-}}
\setlength{\emergencystretch}{2em}
\multlinegap0pt
\usepackage[shortlabels]{enumitem}
\usepackage{xcolor}
\usepackage{amssymb}
\usepackage{tikz}
\usepackage{mathtools}
\usepackage{bm}
\newtheorem{thm}{Theorem}
\DeclareMathOperator{\GL}{GL}
\title{Integrality of $\mathrm{GL}_2\times\mathrm{GL}_2$ Rankin-Selberg integrals for ramified representations}
\author{Alexandros Groutides}
\date{Source: arXiv:2501.16972v2 (CC BY 4.0)}
\begin{document}
\maketitle

\noindent\emph{Source and licence.} Integrality of $\mathrm{GL}_2\times\mathrm{GL}_2$ Rankin-Selberg integrals for ramified representations. Version arXiv:2501.16972v2 (CC BY 4.0), distributed by arXiv under the Creative Commons Attribution 4.0 International licence, http://creativecommons.org/licenses/by/4.0/, as recorded in the arXiv licence metadata for this exact version and shown on the abstract page https://arxiv.org/abs/2501.16972v2. Full licence notice: \texttt{licenses/arxiv-2501.16972-CC-BY-4.0.txt}.\par\smallskip

\noindent\textit{Editorial scope.} Counted unit: the article's thm theorem (source0.tex lines 1484--1495). The article's complete original proof of it is delivered with it (source0.tex lines 1496--1498, one \texttt{proof} environment of 3 lines). No mathematics is written, repaired or re-derived: the statement and the proof are the article's own lines, in the article's own notation, and no retained line is substituted or re-flowed.  No further source-local context is needed: every cross-reference of every retained line resolves inside the two delivered spans.  Nothing was omitted from any retained span, so the package carries no omission record.  No retained line contains a citation, so the wrapper carries no bibliography and no bibliography entry.  No retained line contains a figure environment, an image-inclusion command or drawing code, so no image is delivered and the package declares no asset dependency.  Editorial formatting only, in addition: an AMS \texttt{amsart} preamble that restates the article's own declarations and macros used by the retained lines, namely the environments \texttt{thm} on separate counters, together with the standard packages those lines need.  The retained environments print in this document as Theorem 1 in order of appearance, whereas the article prints the same items with the numbering of its own full document.  The complete native source of the article as distributed in the arXiv e-print for this exact version is bundled unchanged as \texttt{sources/172318/source0.tex}.
\begin{thm}\label{thm: main theorem}
    Let $\pi_1,\pi_2$ be irreducible admissible generic tempered representations of $\GL_2(F)$. Set $\Pi:=\pi_1\times\pi_2$, and let $A\subseteq\mathbf{C}$ be a $\mathbf{Z}[q^{-1},\mu_{\nu q^\tau}]$-algebra containing the spherical Hecke eigenvalues of $\pi_i$ \emph{(}if spherical\emph{)} and the following unramified unitary characters in each case: 
    \begin{enumerate}
        \item $\omega_{\pi_i}$\emph{:} if $\pi_i$ is an unramified principal-series representation.
        \item $\chi_i$\emph{:} if $\pi_i\simeq\mathrm{St}_{\chi_i}$ is an unramified twist of the Steinberg representation
        \item $(\omega_{\pi_i}^\mathrm{unr})^{1/2}$ and $\chi_i$\emph{:} if $\pi_i\simeq (\omega_{\pi_i}^\mathrm{unr})^{1/2}\otimes I(\chi_i\omega,\chi_i^{-1})$ is a half-ramified principal-series representation with $\omega(\varpi)=1$, or $\pi_i\simeq (\omega_{\pi_i}^\mathrm{unr})^{1/2}\otimes I(\chi_{i,1}\chi_i,\chi_{i,2}\chi_i^{-1})$ is a fully-ramified principal-series representation with $\chi_{i,1}(\varpi)=\chi_{i,2}(\varpi)=1$.
        \item $(\omega_{\pi_i}^\mathrm{unr})^{1/2}$\emph{:} if $\pi_i\simeq (\omega_{\pi_i}^\mathrm{unr})^{1/2}\otimes \mathrm{St}_{\chi_i}$ is a ramified twist of the Steinberg representation with $\chi_i(\varpi)=1$, or $\pi_i\simeq (\omega_{\pi_i}^\mathrm{unr})^{1/2}\otimes \pi(\xi_i)$ is a supercuspidal representation with $\omega_{\pi(\xi_i)}(\varpi)=1$.
    \end{enumerate}
    Let $(\phi,g_1,g_2)$ be any $\Pi$-integral datum. Then,  $\Phi(\phi,g_1W_{\pi_1}^\mathrm{new},g_2W_{\pi_2}^\mathrm{new};X)\in A[X,X^{-1}]$ and
    $$Z(\phi,g_1W_{\pi_1}^\mathrm{new},g_2W_{\pi_2}^\mathrm{new};s)=L(\Pi,s)\cdot \Phi(\phi,g_1W_{\pi_1}^\mathrm{new},g_2W_{\pi_2}^\mathrm{new};q^s)$$
    in $L(\Pi,s)A[q^s,q^{-s}]\subseteq A(q^s,q^{-s}).$
\end{thm}

\begin{proof}
    If $\pi_i$ is not an unramified principal-series representation or an unramified twist of the Steinberg representation, then we can re-write it as  $\pi_i\simeq (\omega_{\pi_i}^\mathrm{unr})^{1/2}\otimes \pi_i^{'}$ where $\omega_{\pi_i}^\mathrm{unr}$ is the unramified character defined by $\omega_{\pi_i}^\mathrm{unr}(\varpi):=\omega_{\pi_i}(\varpi)$. The representation $\pi_i^{'}$ is in the same class as $\pi_i$, but with central character that is trivial on $\varpi$, and equal to $\omega_{\pi_i}$ on $\mathcal{O}^\times$. Finally, $c(\pi_i^{'})=c(\pi_i)$. We note that all of our results in this section remain unchanged after unramified twists, and thus we can now apply the required proposition of this section, depending on the classes of $\pi_1,\pi_2$, while potentially enlarging $A$ to ensure it contains $(\omega_{\pi_i}^\mathrm{unr})^{1/2}$. This also implicitly ensures that in every case, the inverse $L$-factor $L(\Pi,s)^{-1}$ is contained in $A[q^s,q^{-s}]$.
\end{proof}

\end{document}
